\section{おわりに}
\label{sec_conclusion}
本稿では，大喜利投稿サイト「ボケて」の約682万件のお題画像に，大喜利の形式，画像中の記号や文字の有無，画像の媒体を表すラベルを，有料の API を用いずに単一の GPU で付与する方法を確立するまでの経緯と分析を報告した．
VLM である Qwen3-VL-8B を教師とし，SigLIP 2 に基づく分類器を生徒とする蒸留は，教師と同等以上の精度を約1桁速い速度で得た．
精度の上限はラベルの定義の曖昧さによって決まっており，ボケの生成に必要な処理から逆算してラベルの定義を4回改訂することで，判定の揃うラベルの体系を得た．
教師の構成では，軸ごとに質問を分け，誤りの多いお題文に限って別の短い質問で確かめる方法が有効であった．
生徒の構成では，生徒の予測で少ない区分を選んで教師ラベルを追加する方法と，埋め込みの MLP と微調整した画像エンコーダの確率を平均する方法が有効であった．
最終の構成は，評価用データで3軸のマクロ F1 の平均 0.919，実データで全項目一致の推定値 0.838 を達成し，全データの予測を約10.2時間で完了した．
全データの集計から，穴埋めがお題の約2割を占めること，2014年前後を境にお題画像が写真中心から文字や作品の画像を含むものへと多様化したことが分かった．

今後の課題として，複数の人の判断者による目視ラベルの一致度の測定，確認の質問を適用した後の全データの無作為抽出による精度の確認，穴埋めの空欄の位置とカルタのひらがなの検出の定量評価が挙げられる．
また，属性の文字ありが大喜利の回答に実際に影響しているかを，ボケの文章を用いて判定する後段の処理の構築も課題である．
